\clearpage
\hypertarget{native-491-lindenau}{}
\label{native-491-lindenau}
\begin{center}
[VON LINDENAU AN GAUSS.]

[In \textsc{Gauss}’ Nachlass.]
\end{center}

\begin{flushright}Altenburg, 17. Febr. 1806.\end{flushright}

--- --- --- Den Brief, in dem Euer Wohlgeboren mir freundlicherweise Ihr Gutachten über \textsc{Rohde}s Erwiderung auf die Kritik mitgeteilt hatten, habe ich Herrn Obh. \textsc{von Zach} weitergegeben. Wie Sie aus der M. C. ersehen haben werden, sah er sich veranlasst, den ganzen Brief abzudrucken. Ich hoffe, das läuft Ihren Wünschen nicht zuwider. --- ---

Ich muss noch einen kleinen Zweifel erwähnen, der mir bei dieser Gelegenheit gekommen ist. Euer Wohlgeb. erklären \textsc{Rohde}s Methode auch in \emph{theoretischer} Hinsicht für falsch, und die Gründe, die Sie dafür anführen, leuchten mir durchaus ein. Bedenke ich andererseits, dass \textsc{Rohde}s Ausdruck, wie ich in meiner Erwiderung gezeigt habe, dem bekannten Verhältnis zwischen Massen, Umlaufzeiten und mittleren Entfernungen analog ist, so ließ mich das wiederum zögern, jene Methode ohne Einschränkung für falsch zu erklären. --- ---

